% Part: computability
% Chapter: tm-computations
% Section: church-turing-thesis

\documentclass[../../../include/open-logic-section]{subfiles}

\begin{document}

\olfileid{tur}{mac}{ctt}
\olsection{Tézis Church--Turing}

Mesin Turing dimaksudake dadi pangganti kang tegesé
cetha kanggo konsep prosedur efektif. Miturut Turing,
sapa wae kang mangerteni konsep prosedur efektif lan
konsep mesin Turing bakal nduwé pangira intuitif yen
apa wae kang bisa ditindakake nganggo prosedur efektif
uga bisa ditindakake dening mesin Turing. Pratelan iki
didhukung dening kasunyatan yen kabeh pangganti
liya kang wis diusulake kanggo menehi teges kang cetha
marang konsep prosedur efektif pranyata padha kanthi
ekstensional karo konsep mesin Turing---yaiku bisa
ngitung persis himpunan fungsi kang padha. Pratelan
iki diarani \emph{tézis Church--Turing}.

\begin{defn}[Tézis Church--Turing]
\emph{Tézis Church--Turing} nyatakake yen apa wae
kang bisa diitung nganggo prosedur efektif uga bisa
diitung dening mesin Turing.
\end{defn}

Tézis Church--Turing digunakake kanthi rong cara.
Cara kapisan dadi dalih kanggo ora gelem repot.
Upamane kita nduwé gambaran prosedur efektif
kanggo ngitung sawijining bab, kayata nganggo
``pseudokode.'' Kita bisa nganggo tézis
Church--Turing kanggo mbenerake pratelan yen ana
mesin Turing kang ngitung fungsi kang padha, senajan
kita durung tenan mbangun mesin iku.

Cara liyane nggunakake tézis Church--Turing luwih
narik kawigaten saka sudut filsafat. Bisa dibuktekake
yen ana fungsi kang ora bisa diitung dening mesin
Turing. Saka kene, kanthi nggunakake tézis
Church--Turing, kita bisa nyimpulake yen fungsi iku
ora bisa diitung kanthi efektif nganggo prosedur
apa wae. Awit yen prosedur kaya mangkono ana,
miturut tézis Church--Turing mesthine uga ana mesin
Turing kanggo fungsi iku. Mula yen kita bisa mbuktekake
yen ora ana mesin Turing kang ngitung fungsi iku,
prosedur efektif uga ora ana. Mligine, tézis
Church--Turing digunakake kanggo nyatakake yen
masalah kang diarani Masalah Mandheg ora mung ora
bisa dirampungake dening mesin Turing, nanging uga
ora bisa dirampungake kanthi efektif babar pisan.

\end{document}
